% =====================================================================
%  ch01.tex —— 第 1 章 序章（原书 p.1–5，PDF 10–14）
% =====================================================================
\chapter{序章}

本章以叙事的方式讲述辛几何，以此对本书所论各个主题作一导引。各主题更精确的内容，将在相应一章的开头详细叙述。

% [原文 p.1]
\section{几何学与物理学}

Platon 在雅典近郊创设的「学园」（Akadēmeia）门上，据说曾悬挂一块匾额，上书「不懂几何者不得入此门」。在希腊时代，几何学在一切学问中占据着特殊的地位。在这个作为理念界之影的世界里，由完全直线与完全圆所构成的几何学，或许正是窥见理念界的一面镜子。

% [Tutor Note: 「アカデミア」即柏拉图于公元前 387 年前后在雅典西北郊创立的学园（Akadēmeia），门匾传说见诸后世记述，真伪不详，但常被引作几何学地位的象征。]

希腊时代最杰出的数学家无疑是 Archimedes。他生于公元前 3 世纪，却远远超越了自己的时代：他事实上完成了静力学，对于积分法也作出了开创性的深刻洞察。此外，在极端轻视实用的文化氛围中，他还是一位能干的工程师。然而，他一生最爱的是几何学。据说他遗言在自己的墓碑上刻下圆柱及其内切球的图形。

% [Tutor Note: 阿基米德墓上刻圆柱与内切球，典出 Cicero《Tusculanae disputationes》；他证明球体积为外切圆柱体积的 2/3，并希望以此作为自己的墓志。]

% [原文 p.2]
然而，他终究未能理解动力学。直到 16 世纪，Galilei 才正确地看出：力与速度无关而与加速度有关；不受力作用的物体作匀速直线运动；在只有重力作用的情形下，重物与轻物同样下落。基于经验积累的直觉被摩擦效应所遮蔽，因此在真空中羽毛与铅球同时落下，对今天的我们而言仍不无惊奇。

然而，尽管只差一步，Galilei 最终还是未能理解运动定律。创立他所欠缺的微积分学、从而完成运动定律的，是当时首屈一指的科学家兼数学家 Newton。他证明了与距离平方成反比的力可以导出关于行星运动的 Kepler 定律。

Newton 对运动定律的表述宣告了近代科学的黎明。其最辉煌的成果，大概是 Halley 等人对周期 75 年的彗星所作的轨道计算。彗星按照 Newton 力学的预言如期回归，由此确认：Newton 力学不仅对地上的物理，对天体的运动也具有强大的预言能力。从地球到太阳，即使以每小时 4 公里、每天 24 小时、一年 365 天不停地走，也要花 4000 年以上；而支配着这种与日常相距极远的天文学尺度上天体运动的定律，与支配我们在地面上抛出的小球轨迹的定律竟是同一个。

Newton 的工作被总结为题为《自然哲学的数学原理》的著作。这本书极具几何色彩。这里的「几何学」是希腊人的几何学，它运用如今已经失传的、关于圆锥曲线的各种结果，因此现代人没有注解便读不下去。

% [Tutor Note: 「ハレー彗星」的回归预言见 Halley 1705 年《Astronomiae cometicae synopsis》；1705 年他据 Newton 力学算出周期约 75 年，并预言 1758 年回归。原书此处以「Halley ら」称之。]

《自然哲学的数学原理》之难解并不限于现代人，对同时代人亦然。然而在数学中，本来一切都应当是自明的；之所以难解，往往是因为用来记述它的语言不够恰当。

Newton 的继承者们背对几何学的世界，致力于建设力学的分析理论。其最初的里程碑是 Lagrange 的分析力学，他在其中高声宣布把一切「图」逐出了力学。离开几何直观、进入公式及其代数运算的世界之后，力学问题的表述与求解便从需要灵感的创造性事业，转变为机械的简单劳动。微积分学与力学成为学生的必修科目，也成为所有工程师的标准工具。

% [原文 p.3]
然而，几何学与物理学之间失去的纽带，在进入 20 世纪后以远为强大的形式复活了。其一是狭义与广义相对论，它们与 Minkowski 几何及伪 Riemann 几何难分难解地联系在一起。此外，以电磁学与量子力学为代表的规范理论，也是用主丛及其联络的语言自然描述、极其几何化的理论。

与此不同的一条脉络是分析力学的几何化。分析力学背后延展着某种几何世界——这一认识自 19 世纪起逐渐成形，进入 20 世纪后变得更加明确。此后 Weyl 发明了「辛（symplectic）」一词，Cartan 则以扩展 Klein 的 Erlangen 纲领的方式创立了 G 结构的几何学。辛流形正是在这一 G 结构几何学的脉络中被自然地定义出来。

% [Tutor Note: 原书作「エルランゲン目録」，指 Klein 1872 年的 Erlangen Programm（埃尔朗根纲领）：以变换群对几何分类。]
% [Tutor Note: 「シンプレクティック」一词由 Weyl 在 1939 年《The Classical Groups》中引入，取自希腊语 sym-plektikos（交织、缠绕）。]

这些脉络历经分岔与汇合而成长，如今已成大江大河，影响着数学与物理学的广泛领域。辛几何在没有明确中心的情况下已成长为庞大的领域，因此难以作总的概观；其主要话题大致可以列举如下。

\takeaway{几何学与力学自 Newton 时代分道扬镳，经相对论、规范理论与分析力学的几何化，在 20 世纪重新合流；辛几何正是这条合流之河的名称。}

\section{群作用}

系统具有对称性时，实质性的自由度减少，问题便容易处理。为自然地处理系统所具有的对称性而建立一般框架，是辛几何发展的动力之一。主张对称性与守恒律本质上是同一件事的 Noether 定理，是其中一个里程碑，也是其后发展的出发点。

Noether 定理具备真正伟大定理的两个特征：主张深远、蕴含众多推论；证明自明。此外，它超出辛几何的框架，在经典力学与量子力学双方都成立，这也是显著特点。在这个意义上，正像同样超越经典论与量子论框架而成立的热力学那样，可以说它标示了人类的一个到达点。

系统的对称性被自然地表述为关于群作用的不变性。当对称性是连续的时候，出现的群是 Lie 群。Lie 群与 Lie 环独立于辛几何而被定义，但辛几何提供了有效运用 Lie 群与 Lie 环概念的绝佳具体例。这在下一项所述的可积系统中尤为显著。反过来，Lie 群的余伴随轨道也给出辛流形的重要例子。

在 Hamilton 系统对称性的脉络中特别重要的，是 Lie 群对辛流形的 Hamilton 作用这一概念。考虑辛流形按 Hamilton 作用取商，便自然地到达辛商的概念；它与代数几何中的几何不变量理论关系极为密切。

\takeaway{对称性即群作用下的不变性；Noether 定理把对称性与守恒律等同起来，而辛几何为 Lie 群理论提供了最自然的舞台。}

\section{可积系统}

任意给定一个 Hamilton 系统时，能具体写出运动方程精确解的情形很少，大多数方程都「解不出来」。于是，方程何时「可解」便自然成为一个问题。

不限于运动方程，对「可解」作精确的表述，自古以来就是数学的重要课题，如今也依然如此。「不可解」问题的代表，大概要数所谓的希腊三大尺规作图问题：三等分角问题（trisecting the angle）、立方倍积问题（doubling the cube）、化圆为方问题（squaring the circle）。这些问题当然用合适的工具就能解决（例如用折纸便可以三等分角、倍积立方体），但若限制所用工具则无法解决。古典上只允许使用圆规与直尺，这些问题自希腊时代起即被研究，直到 19 世纪才终于被证明是不可能的。

% [Tutor Note: 三大作图问题的不可解性分别由 Wantzel（1837，三等分角、倍积立方体）与 Lindemann（1882，化圆为方，依据 $\pi$ 的超越性）证明。]

% [原文 p.4]
「不可解」问题的另一个代表，是关于五次方程求根公式的问题；它经 Ruffini、Abel，由 Galois 以完整的形式解决。方程的 Galois 理论经 Dedekind、Kronecker 等人的贡献，由 Artin 完成为域的 Galois 理论，并被编入大学的标准课程。此后，域的 Galois 理论与覆盖的 Galois 理论相统一，又经 Grothendieck 彻底的抽象化，直至今日。

回到 Hamilton 系统。Hamilton 系统所能具有的可交换对称性，至多是维数的一半；具有达到这一上限的对称性的系统，称为完全可积系统。完全可积性是系统「可解」的一种表述；而描述这种情形下系统结构的，正是 Liouville–Arnold 定理。该定理从三个假设——系统具有维数一半个数的可交换对称性；初值是守恒量的正则点；等值集紧——导出三个结论：等值集是 Lagrange 环面；存在称为作用–角坐标的特殊坐标；系统的运动在作用坐标方向为常量，在角坐标方向为匀速直线运动。

% [Tutor Note: 「等位集合」指所有守恒量取定值的公共水平集；在 Liouville–Arnold 定理的假设下它是 $n$ 维环面 $T^n$，即 Lagrange 环面，见 6.1 节。]

Liouville–Arnold 定理是关于完全可积系统结构的有力且决定性的定理；作为完全可积系统的抽象一般理论，几乎没有再需补充的决定性内容。因此，可积系统的研究必然成为各论。已知的可积系统多种多样，既有孤立的例子，也有在种种背景下系统性地构造出来的例子。其代表是 Toda 格子，它与 Lie 环、Yang–Baxter 方程、Lax 方程等各种话题相关联，颇有趣味。

\takeaway{可积性的上限是「维数一半个数的可交换守恒量」，Liouville–Arnold 定理把这一定义转化为 Lagrange 环面上的作用–角坐标；此后便进入各论。}

\section{动力系统}

可积系统是优美而重要的主题，但给定的 Hamilton 系统可积的情形极为罕见。为了定性地研究不可积系统的性质，Poincaré 创立了动力系统与拓扑学的理论。这一领域的指导性问题是大阳系的稳定性，最终结晶为 KAM（Kolmogorov–Arnold–Moser）理论。

把太阳系的稳定性作为质点系力学的问题来提出，若考虑到以下几点，似乎显得有些奇怪：一旦计入广义相对论效应便有引力波辐射；实际上太阳与行星并非质点，在 Newton 力学中也会因潮汐力而发生耗散；太阳有寿命；而在无宇宙项的广义相对论中，宇宙本身根本不可能定常。尽管如此，这一问题曾吸引了相当多的兴趣，这是历史事实。与此相关，Arnold 提出了 Arnold 猜想：用相空间的拓扑不变量从下方控制周期 Hamilton 系统周期解的个数。这成为此后辛拓扑发展的原动力。

% [原文 p.5]
KAM 理论处理接近可积系统的 Hamilton 系统；而对远离可积系统的 Hamilton 系统发挥威力的，则是遍历理论。在遍历理论中，人们放弃追踪各个初值的时间演化，转而研究系统在长时间中的平均行为。构成遍历理论基石的，是 Liouville 定理。该定理断言 Hamilton 系统的时间演化是关于由辛形式决定的自然测度的保测变换，其证明比 Noether 定理更为自明。

然而，时间演化的保测性是回复定理、遍历定理等重要定理的前提条件，因此 Liouville 定理绝非无聊的定理。毋宁说，Liouville 定理证明的自明性，正是用辛几何描述经典力学的正当性根据之一，也说明我们对经典力学的理解之深。

\takeaway{从 Poincaré 到 KAM 与遍历理论：不可积系统的定性理论把「追踪轨道」换成了「测度与拓扑」，而 Liouville 定理是其自明却不可替代的基石。}

\section{量子化}

本书虽不涉及，但现代的辛几何无法撇开与量子化的关联来谈。辛几何为经典力学提供了优美、精练而又实用的描述框架。另一方面，量子力学给出超越经典力学的、关于自然的更根本的描述，因此在从经典力学过渡到量子力学时，舍弃辛几何、改用全然不同的新框架，似乎也理所当然。

然而不可思议的是，现实并非如此。毋宁说，正是量子化问题使辛几何的重要性与日俱增。要跨越某个领域，就必须把该领域钻研透彻——大概就是这个意思吧。进入 20 世纪后半期，人们又开始探究它与规范理论、超弦理论的关系，辛几何的研究愈发兴盛。

将来，辛几何或许会消亡，人们转而研究「量子几何」，只偶尔作为其经典极限而被提及。但要真的实现这一点、使辛几何只在数学史领域被研究，似乎还相当遥远。

% [Tutor Note: 原书此处作「古典力学から量子力学への以降」，「以降」当为「移行」之误，译文按「过渡」处理。]

\takeaway{量子化非但没有废除辛几何，反而使它的重要性不断增长；辛几何的「消亡」至今仍只是远景。}

\chapterreview{%
本章以物理学史为线索，说明辛几何为何出现：Newton 力学被 Lagrange 改写为分析力学，几何直观退场；20 世纪相对论与规范理论又把几何学带回物理学。由此引出全书的四条主线——群作用与 Noether 定理、完全可积系统与 Liouville–Arnold 定理、不可积系统的动力系统与 KAM/遍历理论、以及量子化。
原书强调，辛几何的价值不仅在「描述」，更在于它把对称性、守恒律、可积性与拓扑不变量编织成同一套语言。
}
